\documentclass[11pt,a4paper]{article}
% Preámbulo compartido para los modelos del segundo parcial de ACN.
% Replica el formato del parcial real 2P-2S-2025.
\usepackage[utf8]{inputenc}
\usepackage[T1]{fontenc}
\usepackage[spanish]{babel}
\usepackage{amsmath,amssymb}
\usepackage[margin=2.6cm]{geometry}
\usepackage{enumitem}
\usepackage{array}

\pagestyle{plain}
\setlength{\parindent}{0pt}
\setlength{\parskip}{5pt}

% Listas de incisos: a), b), c), ...
\setlist[enumerate,1]{label=\alph*),leftmargin=2.2em,itemsep=7pt,topsep=7pt}
% Listas numeradas 1., 2., 3., ...
\newlist{puntos}{enumerate}{1}
\setlist[puntos]{label=\arabic*.,leftmargin=2.2em,itemsep=7pt,topsep=7pt}
% Listas de condiciones del enunciado
\setlist[itemize,1]{label=$\bullet$,leftmargin=2em,itemsep=5pt,topsep=6pt}

% Encabezado del parcial
\newcommand{\encabezado}[1]{%
  \begin{center}
    \textbf{\large Aplicaciones Computacionales en Negocios}\\[3pt]
    \textbf{\large Segundo semestre de 2025 --- Segundo parcial}\\[3pt]
    \textbf{\normalsize Modelo de práctica #1}
  \end{center}
  \vspace{4pt}
  \hrule height 0.8pt
  \vspace{14pt}
  Nombre:\par\vspace{12pt}
  Legajo:\par\vspace{12pt}
  DNI:\par\vspace{12pt}
  ¿Te recibís con este parcial?:
  \vspace{10pt}
  \hrule height 0.8pt
  \vspace{16pt}
}

% Título de ejercicio
\newcommand{\ejercicio}[1]{\par\medskip\noindent\textbf{Ejercicio #1.}\ }

% Atajos de notación usados por la cátedra
\newcommand{\R}{\mathbb{R}}
\newcommand{\Z}{\mathbb{Z}}
\newcommand{\dem}{\mathrm{dem}}
\newcommand{\capa}{\mathrm{cap}}
\newcommand{\costo}{\mathrm{costo}}


\begin{document}
\encabezado{1}

%======================================================================
\ejercicio{1}
Consideremos el modelo clásico de \emph{lot sizing} con un único producto, con
costos de producción y de almacenamiento. Tenemos un conjunto
$T = \{1,\dots,n\}$ de períodos de planificación; para cada $t \in T$ tenemos la
demanda $d_t \in \R_+$, el costo unitario de producción $c_t \in \R_+$ y el costo
unitario de almacenamiento $a_t \in \R_+$. La variable $x_t \in \R_+$ representa
la cantidad fabricada en el período $t$, y la variable $s_t \in \R_+$ representa
el stock al final del período $t$:

\begin{align}
\min \quad & \sum_{t \in T} \big( c_t\, x_t + a_t\, s_t \big) \\
s_t &= s_{t-1} + x_t - d_t && \forall\, t \in T \\
x_t &\ge 0 && \forall\, t \in T \\
s_t &\ge 0 && \forall\, t \in T
\end{align}

con $s_0 = 0$ (el stock inicial es nulo).

\begin{enumerate}
  \item La restricción (2), junto con la restricción (4), es la que impone el
  cumplimiento de la demanda en cada período. Explicar por qué. ¿Qué
  representaría operativamente el modelo si reemplazáramos (4) por
  $s_t \in \R$ (es decir, si permitiéramos stock negativo)? En ese caso,
  ¿cómo habría que modificar la función objetivo para que el modelo siga
  teniendo sentido?

  \item Supongamos ahora que la planta tiene una capacidad máxima de producción
  $C_t$ en cada período, y que encender la línea de producción en un período
  tiene un costo fijo $f_t \in \R_+$ (que es dato), independiente de la cantidad
  fabricada. ¿Cómo se modifica el modelo? Supongamos además que, si en un
  período se produce, el lote fabricado debe ser de al menos $m$ unidades
  (que también es dato). ¿Cómo se agrega esta condición?

  \item Por razones de mantenimiento, no se puede producir durante tres períodos
  consecutivos. Agregar esta restricción al modelo del punto anterior.
  Supongamos además que entre dos períodos con producción debe haber al menos
  $\ell$ períodos sin producción. ¿Cómo se escribe esta restricción?

  \item ¿Cómo se modifica el modelo original si el objetivo no es minimizar el
  costo total, sino minimizar el costo más elevado entre los períodos de $T$?
  Por otro lado, volviendo al modelo original: ¿qué sucedería con las soluciones
  óptimas si para algún período tuviéramos $a_t < 0$? ¿Y si además el depósito
  no tuviera capacidad máxima?
\end{enumerate}

\newpage
%======================================================================
\ejercicio{2}
Estamos a cargo de la planificación de las guardias de un hospital durante las
próximas cuatro semanas. Sea $M$ el conjunto de médicos y sea $E$ el conjunto de
enfermeros, y llamamos $T = \{1,\dots,28\}$ a los días del período de
planificación. Cada día está dividido en dos turnos: el turno \emph{día} y el
turno \emph{noche}, y llamamos $S = \{\mathrm{d}, \mathrm{n}\}$ al conjunto de
turnos. Tenemos que determinar qué turnos deberá cubrir cada empleado,
respetando las siguientes restricciones:

\begin{itemize}
  \item Cada empleado trabaja a lo sumo un turno por día.
  \item Cada empleado debe cubrir entre 16 y 18 turnos a lo largo de las cuatro
  semanas.
  \item Ningún empleado puede trabajar el turno noche de un día y el turno día
  del día siguiente.
  \item Ningún empleado puede trabajar más de cuatro días consecutivos
  (sin importar el turno que cubra en cada uno de esos días).
  \item En cada turno de cada día debe haber al menos dos médicos y al menos
  tres enfermeros.
  \item En ningún turno puede haber más de ocho empleados trabajando.
\end{itemize}

\begin{puntos}
  \item Formular un modelo de programación lineal entera para encontrar una
  solución a este problema, buscando minimizar la cantidad máxima de turnos
  noche que recibe un mismo empleado. \emph{Pista:} para cada
  $i \in M \cup E$, cada $k \in S$ y cada $t \in T$, considerar una variable
  binaria $x_{ikt}$ que represente si el empleado $i$ cubre el turno $k$ del
  día $t$. Observar que la función objetivo pedida no es lineal: explicar cómo
  se linealiza.

  \item Sea $W \subseteq T$ el conjunto de los días que caen sábado o domingo,
  y supongamos que los días $t$ y $t+1$ forman un fin de semana cuando
  $t \in W$ y $t+1 \in W$. Agregar al modelo una restricción para que cada
  médico tenga al menos un fin de semana completo libre.

  \item Supongamos ahora que el hospital quiere además que ningún empleado
  cubra más de tres turnos noche en una misma semana. Agregar esta restricción
  al modelo. ¿Puede este agregado volver infactible el problema? Justificar.
\end{puntos}

\newpage
%======================================================================
\ejercicio{3}
Estamos a cargo de la organización de una campaña de testeo sanitario y debemos
determinar qué centros de atención abrir y cuántos \emph{boxes} habilitar en cada
uno. Tenemos un conjunto $L$ de ubicaciones candidatas y un conjunto $C$ de
barrios. Para cada barrio $i \in C$ conocemos su demanda $\dem_i \in \Z_+$
(cantidad de personas a testear) y, para cada ubicación $j \in L$, la distancia
$d_{ij} \in \R_+$ entre el barrio $i$ y la ubicación $j$. Todas las personas de un
mismo barrio deben ser atendidas en el mismo centro.

Para cada ubicación $j \in L$ tenemos la cantidad mínima $a_j \in \Z_+$
(en caso de apertura) y la cantidad máxima $b_j \in \Z_+$ de boxes que se pueden
habilitar allí; es decir, la ubicación $j$ puede tener 0 boxes (si no se abre) o
bien cualquier número de boxes entre $a_j$ y $b_j$. Cada box atiende hasta 120
personas. Para cada ubicación $j \in L$ tenemos además el costo $c_j \in \R_+$ de
habilitación de cada box y un costo fijo $f_j \in \R_+$ por abrir el centro (es
decir, si se habilitan $k$ boxes en la ubicación $j$ entonces se incurre en un
costo de $f_j + k\,c_j$, y si no se habilita ningún box en la ubicación $j$
entonces el costo asociado a esa ubicación es cero).

\begin{puntos}
  \item Formular un modelo de programación lineal entera para determinar qué
  centros abrir, cuántos boxes habilitar en cada uno y en qué centro se atiende
  cada barrio, de modo tal que toda la demanda quede cubierta y se minimice el
  costo total. Observar que en este punto no es relevante minimizar la distancia
  de viaje de las personas.

  \item Agregar al modelo una restricción para que ningún barrio sea asignado a
  un centro ubicado a más de $U$ kilómetros de distancia (siendo $U$ un dato).
  ¿Puede este agregado volver infactible el problema? ¿Qué se suele hacer en
  esos casos?

  \item Supongamos ahora que disponemos de un presupuesto total $B$ y que, en
  lugar de minimizar el costo, queremos minimizar la distancia que debe recorrer
  el barrio que queda más lejos de su centro asignado. ¿Cómo se modifica el
  modelo? Explicar la linealización utilizada.

  \item Supongamos finalmente que cada barrio $i \in C$ tiene una distancia
  actual $\mathrm{actual}_i \in \R_+$ al centro donde se atiende hoy, y que
  queremos garantizar que ningún barrio empeore su situación respecto de la
  actual. ¿Cómo se agrega esta condición al modelo? ¿Por qué conviene
  implementarla fijando variables en lugar de agregar desigualdades nuevas?
\end{puntos}

\end{document}
